\begin{proof}[Proof of \cref{obj:prop:unrestricted-fixed-q-phase-boundaries}]
Fix \(q\), the arrival floor, with \(0<q<1/2\). Throughout the proof, \(n\) ranges over positive integers, and all limits are taken as \(n\to\infty\).

\begin{enumerate}
\item For a sequence of integers \(d_n\ge1\), define
\[
\ell_n^{\mathrm{fix}}:=\log(e+n),
\qquad
r_n^{\mathrm{fix}}
:=
\min\left\{1,\frac1n+
\left(\frac{d_n}{n\ell_n^{\mathrm{fix}}}\right)^2\right\}.
\]
Apply \cref{obj:thm:unrestricted-fixed-q-minimax-frontier}, whose lower comparison uses the published zero-control-regression lower bound
\citep[Theorem~2, Section~3.2, p.~9, and its proof in Section~C.5, pp.~39--44]{ZengBalakrishnanHanKennedy2026}.
It supplies constants \(0<\underline c(q)\le\overline C(q)\), depending only on \(q\), such that
\[
\underline c(q)\,r_n^{\mathrm{fix}}
\le \mathfrak R^{+}_{n,d_n,q}
\le \overline C(q)\,r_n^{\mathrm{fix}}
\qquad(n\ge1).
\]
These constants are common to every dimension sequence. Moreover,
\[
0\le r_n^{\mathrm{fix}},
\qquad
0\le\overline C(q).
\]
% lean: CausalSmith.Stat.MarRareqLogfrontier.unrestricted_fixed_q_minimax_frontier

\item We first establish the consistency equivalence. The comparison above gives
\[
0\le r_n^{\mathrm{fix}}
\le \underline c(q)^{-1}\mathfrak R^{+}_{n,d_n,q},
\qquad
0\le\mathfrak R^{+}_{n,d_n,q}
\le\overline C(q)\,r_n^{\mathrm{fix}}.
\]
Applying the squeeze theorem in each direction yields
\[
\mathfrak R^{+}_{n,d_n,q}\longrightarrow0
\quad\Longleftrightarrow\quad
r_n^{\mathrm{fix}}\longrightarrow0.
\]

To characterize the latter limit, define
\[
\nu_n:=\frac1n,
\qquad
\xi_n:=\frac{d_n}{n\ell_n^{\mathrm{fix}}}.
\]
Since \(e+n>1\), we have
\[
\ell_n^{\mathrm{fix}}>0,
\qquad
n\ell_n^{\mathrm{fix}}>0,
\qquad
\xi_n\ge0,
\qquad
\nu_n\ge0,
\qquad
\nu_n\longrightarrow0.
\]
The last limit follows from \(n\to\infty\) and continuity of inversion at infinity. Continuity of the square root and nonnegativity of \(\xi_n\) give
\[
\xi_n^2\longrightarrow0
\quad\Longrightarrow\quad
\xi_n=\sqrt{\xi_n^2}\longrightarrow0.
\]
Conversely, continuity of squaring gives
\[
\xi_n\longrightarrow0
\quad\Longrightarrow\quad
\xi_n^2\longrightarrow0.
\]
The positive denominator also gives the usual quotient characterization of little-\(o\):
\[
d_n=o\!\left(n\ell_n^{\mathrm{fix}}\right)
\quad\Longleftrightarrow\quad
\xi_n\longrightarrow0.
\]

It remains to handle the truncation. If
\(r_n^{\mathrm{fix}}\to0\), then eventually
\[
r_n^{\mathrm{fix}}=\min\{1,\nu_n+\xi_n^2\}<1.
\]
At every such index,
\[
\nu_n+\xi_n^2<1,
\qquad
r_n^{\mathrm{fix}}=\nu_n+\xi_n^2:
\]
indeed, \(\nu_n+\xi_n^2\ge1\) would make the minimum equal to \(1\).
Thus \(\nu_n+\xi_n^2\to0\), and
\[
0\le\xi_n^2\le\nu_n+\xi_n^2
\]
implies \(\xi_n^2\to0\). Conversely, if \(\xi_n^2\to0\), then
\[
\nu_n+\xi_n^2\longrightarrow0,
\qquad
\min\{1,\nu_n+\xi_n^2\}\longrightarrow\min\{1,0\}=0,
\]
by continuity of addition and the minimum.
% lean: CausalSmith.Stat.MarRareqLogfrontier.phase_min_rate_tendsto

Consequently,
\[
r_n^{\mathrm{fix}}\longrightarrow0
\quad\Longleftrightarrow\quad
\xi_n^2\longrightarrow0
\quad\Longleftrightarrow\quad
\xi_n\longrightarrow0
\quad\Longleftrightarrow\quad
d_n=o\!\left(n\log(e+n)\right).
\]
Together with the risk comparison, this proves the first assertion.
% lean: CausalSmith.Stat.MarRareqLogfrontier.fixedQRate_tendsto_iff_dimension_littleO

\item For the sufficient condition for the parametric rate, choose
\[
c:=\underline c(q)>0.
\]
Given \(M>0\), define
\[
C:=\overline C(q)(1+M^2).
\]
These constants satisfy
\[
c=\underline c(q)
\le\overline C(q)
\le\overline C(q)(1+M^2)=C.
\]
Suppose that eventually
\[
d_n\le M\sqrt n\,\ell_n^{\mathrm{fix}}.
\]
Both sides are nonnegative, so squaring this inequality and using
\((\sqrt n)^2=n\) gives
\[
d_n^2
\le\left(M\sqrt n\,\ell_n^{\mathrm{fix}}\right)^2
=M^2n\left(\ell_n^{\mathrm{fix}}\right)^2.
\]
Dividing by the positive squared denominator yields
\[
\xi_n^2
=\frac{d_n^2}{(n\ell_n^{\mathrm{fix}})^2}
\le\frac{M^2}{n}.
\]
Here the cancellation uses
\[
\frac1n(n\ell_n^{\mathrm{fix}})^2
=n\left(\ell_n^{\mathrm{fix}}\right)^2.
\]
Since \(1/n\le1\) for \(n\ge1\), both entries of the minimum defining the rate are at least \(1/n\). Hence
\[
\frac1n
\le r_n^{\mathrm{fix}}
=\min\left\{1,\frac1n+\xi_n^2\right\}
\le\frac1n+\xi_n^2
\le\frac{1+M^2}{n}.
\]
Multiplying the lower and upper bounds by the corresponding frontier constants gives, eventually,
\[
\frac{c}{n}
\le\underline c(q)\,r_n^{\mathrm{fix}}
\le\mathfrak R^{+}_{n,d_n,q}
\le\overline C(q)\,r_n^{\mathrm{fix}}
\le\frac{C}{n}.
\]
The choice of \(c\) depends only on \(q\), and the choice of \(C\) depends only on \(q\) and \(M\). The eventual dimension bound determines the threshold along each sequence.
% lean: CausalSmith.Stat.MarRareqLogfrontier.unrestricted_fixed_q_phase_boundaries

\item For the converse, suppose there exist constants \(0<c_*\le C_*\) such that eventually
\[
\frac{c_*}{n}
\le\mathfrak R^{+}_{n,d_n,q}
\le\frac{C_*}{n}.
\]
In particular, \(C_*\ge0\). Combining the upper risk bound with the lower frontier comparison gives
\[
\underline c(q)\,r_n^{\mathrm{fix}}
\le\mathfrak R^{+}_{n,d_n,q}
\le\frac{C_*}{n},
\]
and therefore
\[
0\le r_n^{\mathrm{fix}}
\le\frac{C_*}{\underline c(q)}\,\frac1n
\longrightarrow0.
\]
Thus eventually \(r_n^{\mathrm{fix}}<1\).

Define the dimension-bound constant
\[
K_{\mathrm{dim}}
:=\max\left\{1,\frac{C_*}{\underline c(q)}\right\}.
\]
Its calibration gives
\[
0\le K_{\mathrm{dim}},
\qquad
1\le K_{\mathrm{dim}},
\qquad
\frac{C_*}{\underline c(q)}
\le K_{\mathrm{dim}}
\le K_{\mathrm{dim}}^2.
\]
At all sufficiently large indices, \(r_n^{\mathrm{fix}}<1\) implies
\[
\frac1n+\xi_n^2<1,
\qquad
r_n^{\mathrm{fix}}=\frac1n+\xi_n^2,
\]
because a sum at least \(1\) would make the truncated rate equal to \(1\).
The preceding rate bound and \(1/n\ge0\) now yield
\[
\xi_n^2
\le\frac1n+\xi_n^2
=r_n^{\mathrm{fix}}
\le\frac{C_*}{\underline c(q)}\,\frac1n.
\]
Multiplying by \((n\ell_n^{\mathrm{fix}})^2>0\) gives
\[
d_n^2
\le\frac{C_*}{\underline c(q)}
\,n\left(\ell_n^{\mathrm{fix}}\right)^2
\le K_{\mathrm{dim}}^2
\,n\left(\ell_n^{\mathrm{fix}}\right)^2.
\]
The second inequality uses
\(n(\ell_n^{\mathrm{fix}})^2\ge0\). Also,
\[
K_{\mathrm{dim}}\sqrt n\,\ell_n^{\mathrm{fix}}\ge0,
\qquad
\left(K_{\mathrm{dim}}\sqrt n\,\ell_n^{\mathrm{fix}}\right)^2
=K_{\mathrm{dim}}^2n\left(\ell_n^{\mathrm{fix}}\right)^2.
\]
Since \(d_n\ge0\), comparison of these squares implies
\[
d_n\le K_{\mathrm{dim}}\sqrt n\,\ell_n^{\mathrm{fix}}
\qquad\text{eventually}.
\]
Both the dimension and the comparison scale are nonnegative, so this is precisely the required bound
\[
d_n=O\!\left(\sqrt n\,\log(e+n)\right).
\]
% lean: CausalSmith.Stat.MarRareqLogfrontier.unrestricted_fixed_q_phase_boundaries
\end{enumerate}
\end{proof}
